% Chapter 2

\chapter{Literature review} % Main chapter title
\label{chap:Chapter2} % For referencing the chapter elsewhere, use \ref{chap:Chapter2} 

%----------------------------------------------------------------------------------------
\section{PRISMA Methodology Overview}
A systematic literature review will be conducted following \ac{PRISMA} guidelines to ensure comprehensive and reproducible coverage of the research domain. The search strategy encompasses the following elements:

\textbf{Information sources:}
\begin{itemize}
    \item \textbf{Academic databases:} IEEE Xplore, ACM Digital Library, Springer, ScienceDirect, and Google Scholar for peer‑reviewed publications
    \item \textbf{Technical repositories:} arXiv for preprints and working papers in systems and software engineering
    \item \textbf{Gray literature:} Technical documentation from established open‑source observability projects, industry whitepapers, and practitioner reports from reputable sources (given the rapid evolution of container observability practices)
\end{itemize}

\textbf{Search terms and queries:}
Core search queries will combine terms from three concept groups using Boolean operators:
\begin{enumerate}
    \item \textit{Observability concepts:} ``observability'', ``monitoring'', ``telemetry'', ``instrumentation'', ``metrics'', ``logging'', ``tracing'', ``distributed tracing''
    \item \textit{Container and deployment context:} ``container'', ``containerized'', ``Docker'', ``non‑privileged'', ``unprivileged'', ``sidecar'', ``agent''
    \item \textit{Application domain:} ``CI/CD'', ``continuous integration'', ``continuous deployment'', ``DevOps'', ``reliability'', ``incident response'', ``MTTD'', ``MTTR''
\end{enumerate}

Example query structure: (\textit{observability} OR \textit{monitoring} OR \textit{telemetry}) AND (\textit{container} OR \textit{containerized} OR \textit{non‑privileged}) AND (\textit{CI/CD} OR \textit{reliability} OR \textit{incident response})

\textbf{Inclusion and exclusion criteria:}
\begin{itemize}
    \item \textbf{Inclusion:} Peer‑reviewed publications from 2018–2025, technical reports and documentation addressing container telemetry collection, instrumentation patterns, performance evaluation, or reliability improvement in constrained or heterogeneous environments, comparative studies of observability approaches, and empirical evaluations of monitoring overhead
    \item \textbf{Exclusion:} Publications exclusively focused on privileged host access patterns without discussion of constraints, cloud‑provider‑specific solutions that cannot generalize to on‑premise infrastructure, purely theoretical work without practical implementation guidance, and papers predating 2018 (unless seminal works establishing foundational concepts)
\end{itemize}

\textbf{Selection process:}
Literature selection will follow a three‑stage process: (1) title and abstract screening to identify potentially relevant sources, (2) full‑text review to assess relevance against inclusion/exclusion criteria, and (3) quality assessment focusing on methodological rigor, reproducibility, and applicability to constrained environments. Reference chaining (backward and forward citation tracking) will supplement database searches to identify additional relevant work.
\subsection{Identification}
The identification stage will record search queries, databases queried, and the number of retrieved records. Queries will be versioned and documented to enable replication of the search process.

\subsection{Screening}
Title and abstract screening will be performed using documented inclusion/exclusion criteria; screening decisions and reasons for exclusion at this stage will be recorded. When uncertain, items will be advanced to full\-text review to avoid premature exclusion.

\subsection{Eligibility}
Full\-text review assesses whether the study meets the inclusion criteria and whether it reports sufficient methodological detail for quality assessment. Studies that lack sufficient detail will be noted but may be excluded from quantitative synthesis.

\subsection{Inclusion and Data Extraction}
Included studies will have data extracted into structured evidence tables. Extraction fields include: study context, deployment environment (cloud/on\-premise), privilege assumptions, telemetry modalities used, evaluation metrics and experimental design. Extraction will be double\-checked for accuracy.
\section{Prisma Flow Diagram}
The PRISMA flow diagram will be generated from the recorded counts at each stage (identification, screening, eligibility, inclusion) and included in the thesis as a visual summary of the search process.

\subsection{Quality assessment criteria}
Quality assessment will evaluate internal validity (study design, controls), external validity (generalizability to constrained \ac{CICD} environments), and reporting transparency (availability of artifacts, measurement procedures). Assessment outcomes will guide whether a study contributes to quantitative synthesis or only qualitative insight.
\section{Related Work on Observability in CICD Environments}

The literature review will examine existing approaches to container observability, deployment patterns for telemetry collection (agent‑based, sidecar, host‑level), instrumentation frameworks and their privilege requirements, performance and overhead studies, and empirical evaluations of monitoring interventions.

Preliminary reading indicates that much of the documented practice and empirical evaluation in the literature assumes elevated privileges, homogeneous cloud environments and a high degree of control over service code and deployment (for example, the use of sidecar patterns and Kubernetes operators). There is comparatively less empirical guidance for internal, on\-premise \ac{CICD} deployments where host administration is restricted and third\-party packaged services cannot be modified.

Below we summarise representative studies and identify areas where evidence is limited or absent. These summaries are intended to guide the systematic review rather than to be exhaustive.
%----------------------------------------------------------------------------------------

\section{Answering the Research Question}
The systematic review will produce an evidence base to answer the research questions by: (a) identifying feasible telemetry collection patterns under privilege constraints; (b) documenting evaluation methods and KPIs suitable for constrained environments; and (c) providing references and examples to inform the prototype implementation and measurement strategy.

%----------------------------------------------------------------------------------------

\section{Instrumentation and Deployment Assumptions in the Literature}
\label{sec:instrumentation-assumptions}

The comparative analysis of published observability systems reveals a consistent pattern of deployment and privilege assumptions that influence the applicability of proposed solutions to constrained environments. A substantial proportion of evaluated architectures presume host-level administrative access, either through privileged containers with direct socket access, kernel-level instrumentation frameworks, or cluster-wide operator deployments. Correira et al.~\cite{correiaMonintainerOrchestrationindependentExtensible2023} exemplify this pattern with a hybrid push-pull architecture requiring node-level agents with \texttt{cgroup} access for container metrics and RAPL interfaces for energy profiling, capabilities unavailable without elevated privileges. Similarly, Dinga et al.~\cite{dingaEmpiricalEvaluationEnergy2023} deploy host-based collectors such as Metricbeat and cAdvisor with direct Docker API access, supplemented by application-level Zipkin tracing that necessitates explicit integration with containerized services. These deployment patterns, while effective in homogeneous cloud platforms where orchestrators manage infrastructure, present significant barriers in heterogeneous or restricted environments where node access is administratively controlled.

A second class of instrumentation approaches leverages kernel observability mechanisms, particularly \ac{eBPF}, to achieve non-intrusive telemetry collection. Works such as Sharma et al.~\cite{sharmaEBPFEnhancedCompleteObservability2024} and Kannan et al.~\cite{kannanDesigningLightweightNetwork2024} demonstrate that kernel-level probes can extract rich observability signals—including network flow metrics, system call traces, and container-level resource utilization—without requiring application code modifications. However, these mechanisms impose their own set of deployment prerequisites: Linux kernel versions that support stable \ac{eBPF} interfaces (typically v5.15 or newer), capabilities such as \texttt{CAP\_BPF} and \texttt{CAP\_NET\_ADMIN}, and deployment as privileged DaemonSets with node-wide scope. While such techniques offer substantial reductions in instrumentation overhead compared to sidecar proxies or per-service agents, they remain inapplicable in scenarios where host kernel access is restricted or where execution occurs on heterogeneous platforms without uniform kernel feature availability. The distinction between operational feasibility and privilege requirements emerges as a critical dimension in assessing the transferability of observability architectures to environments lacking administrative control over underlying infrastructure.

Sidecar-based collection represents a third architectural family, evaluated extensively in service mesh contexts by Hasanth et al.~\cite{hasanthEvaluatingPerformanceSidecarbased2024} and Salcedo-Navarro et al.~\cite{salcedo-navarroK8sidecarModularKubernetes2025}. Sidecars intercept network traffic at the pod boundary, enabling transparent telemetry collection without modifying application logic. Yet sidecar deployments introduce resource multipliers proportional to service count—each sidecar consumes CPU, memory, and network bandwidth—and rely on orchestrator-managed injection mechanisms that may not be available or permissible in restrictive \ac{CICD} environments. Furthermore, latency analysis by Sahu et al.~\cite{sahuSidecarsCentralLane2023} indicates measurable microarchitectural contention between application and sidecar containers sharing vCPU resources, underscoring that deployment simplicity does not eliminate performance implications.

%----------------------------------------------------------------------------------------

\section{Performance Evaluation Practices and Reported Overheads}
\label{sec:evaluation-practices}

Empirical evaluations in the reviewed literature exhibit considerable heterogeneity in experimental design, workload selection, and performance metrics, complicating direct comparisons across studies. Benchmark applications vary from synthetic microservice suites such as Online Boutique~\cite{hasanthEvaluatingPerformanceSidecarbased2024} and DeathStarBench~\cite{TraceWeaverDistributedRequest} to production-scale workloads deployed in commercial platforms~\cite{liuJCallGraphTracingMicroservices2019}. Load generation techniques range from controlled concurrency ramps using tools such as Locust to stochastic arrival models with diverse traffic patterns. Overhead is quantified through disparate metrics—CPU utilization percentages, absolute memory footprints, throughput degradation relative to baseline, and latency percentiles—reported at different granularities (per-node, per-container, system-wide aggregates). This methodological diversity reflects the absence of standardized observability benchmarking protocols for containerized systems, a gap that hinders reproducibility and cross-study synthesis.

Reported overhead magnitudes span a wide range depending on instrumentation technique and signal fidelity. Host-based \ac{eBPF} solutions demonstrate notably lower resource consumption than alternative approaches: Sharma et al.~\cite{sharmaEBPFEnhancedCompleteObservability2024} report CPU usage reductions of twenty-one to two-hundred-fourteen times compared to OpenTelemetry agents, sidecar proxies, and traditional host exporters. Kannan et al.~\cite{kannanDesigningLightweightNetwork2024} document approximately ten percent throughput degradation and point-five percent per-node CPU overhead for network flow monitoring via Traffic Control hooks, contrasting favorably with legacy packet capture tools that impose fifty to seventy percent throughput penalties. Conversely, full-stack distributed tracing through export-based OpenTelemetry collectors incurs substantial costs: performance evaluations on serverless and microservice workloads reveal throughput reductions ranging from nineteen to eighty percent and latency increases between seven and one-hundred-seventy-five percent~\cite{InvestigatingPerformanceOverhead}. Tail-based sampling strategies, while improving trace completeness, amplify collector CPU consumption by seventy-one percent relative to head-based sampling~\cite{karkanPerformanceOverheadOpenTelemetry}. These figures underscore the existence of fundamental tradeoffs between signal completeness, collection latency, and resource overhead, tradeoffs that must be navigated in resource-constrained deployment contexts.

Several studies acknowledge methodological limitations that restrict the generalizability of overhead assessments. Dinga et al.~\cite{dingaEmpiricalEvaluationEnergy2023} note the inability to isolate container-level energy consumption using external power measurement, attributing observed increases to host system-wide effects rather than application-specific telemetry costs. Norgren~\cite{norgrenOPTIMIZINGDISTRIBUTEDTRACING} highlights that collector placement decisions—DaemonSet versus sidecar—interact with orchestrator resource limits and TraceContext propagation overhead in ways that are not easily predicted from component-level microbenchmarks. Ramachandran et al.~\cite{ramachandranFUSEFaultDiagnosis2023} document latency amplifications between thirteen and forty-five times baseline when employing comprehensive \ac{eBPF} syscall tracing for fault diagnosis, demonstrating that even low-overhead kernel mechanisms impose non-negligible penalties when probe invocation frequency is high. The variability of reported overheads reflects not only differences in experimental setup but also the sensitivity of performance outcomes to workload characteristics, instrumentation scope, and backend processing efficiency.

%----------------------------------------------------------------------------------------

\section{Discrepancies Between Research Environments and Constrained CI/CD Contexts}
\label{sec:environment-mismatch}

The operational assumptions underlying much of the published research diverge substantially from the constraints characteristic of on-premise \ac{CICD} deployments in heterogeneous or administratively restricted environments. The majority of evaluated systems presume orchestrator-mediated infrastructure management, often Kubernetes, with the authority to deploy cluster-wide resources such as DaemonSets, CustomResourceDefinitions, and admission controllers. This model affords centralized configuration, uniform telemetry pipelines, and the ability to inject sidecars or agents into workloads declaratively. By contrast, \ac{CICD} systems operating on shared infrastructure with limited user privileges cannot rely on such orchestrator capabilities. Executors may run as unprivileged containers without access to host namespaces, kernel interfaces, or Docker sockets, precluding the direct application of host-based or kernel-level instrumentation techniques documented in the literature.

Furthermore, research deployments typically assume homogeneous execution environments—identical kernel versions, uniform container runtimes, consistent network plugins—facilitating predictable instrumentation behavior. Real-world \ac{CICD} infrastructure, particularly in organizations with diverse legacy systems or multi-cloud deployments, exhibits platform heterogeneity: varying operating system distributions, mixed kernel versions, and differing container runtime policies. Observability mechanisms dependent on specific kernel features, such as \ac{eBPF} tracepoints introduced in recent Linux releases, fail silently or degrade gracefully depending on fallback support, introducing operational fragility. The assumption of code modifiability presents another divergence: academic prototypes often demonstrate instrumentation by integrating tracing libraries into purpose-built microservice applications, whereas \ac{CICD} pipelines frequently execute third-party packaged software—build tools, testing frameworks, dependency resolvers—that cannot be recompiled or instrumented without violating support contracts or introducing version drift.

A third area of misalignment concerns the evaluation of correlation mechanisms across observability signals. Research systems implementing distributed tracing commonly propagate trace context through HTTP headers or gRPC metadata, mechanisms that require either application-level instrumentation or transparent proxy interception at service boundaries. In \ac{CICD} contexts, correlation must often be reconstructed post hoc from logs, metrics, and execution metadata available through pipeline orchestration APIs, without the benefit of in-band trace propagation. This retrospective correlation problem—linking container resource spikes to specific job identifiers, associating log error patterns with pipeline failure events—remains underexplored in the literature, which tends to assume instrumentation frameworks designed for long-lived service deployments rather than ephemeral task-based workloads. The temporal and architectural characteristics of \ac{CICD} executions—short-lived containers, stateless task execution, batch job semantics—differ meaningfully from the microservice interaction patterns that dominate the observability research corpus.

%----------------------------------------------------------------------------------------

\section{Practitioner-Oriented Approaches and Industry Practices}
\label{sec:practitioner-approaches}

Beyond the academic literature, industry documentation and technical reports from observability platform vendors provide complementary perspectives on telemetry collection in production environments. These sources emphasize operational pragmatism, configuration flexibility, and the incremental adoption of observability capabilities within existing infrastructure constraints. For instance, OpenTelemetry project documentation outlines collector deployment patterns ranging from agent-per-host configurations to centralized gateway architectures, explicitly addressing tradeoffs between local buffering, network bandwidth consumption, and backend ingestion capacity. Prometheus operational guidelines recommend pull-based metric collection with service discovery mechanisms that decouple telemetry exporters from backend topology, enabling declarative configuration suitable for dynamic containerized environments. Grafana Loki documentation describes log aggregation without indexing full message content, reducing storage overhead at the cost of query expressiveness, a tradeoff relevant to resource-constrained deployments.

Vendor technical whitepapers and case studies, while less methodologically rigorous than peer-reviewed research, document deployment experiences across diverse organizational contexts. These reports frequently highlight challenges not emphasized in academic work: authentication and authorization models for multi-tenant telemetry pipelines, data retention policies balancing compliance requirements against storage costs, and observability tool sprawl arising from heterogeneous technology stacks. Datadog and Dynatrace technical blogs discuss the operational complexity of maintaining unified observability across hybrid cloud and on-premise infrastructure, noting that privilege limitations in customer-managed environments necessitate agent deployment models that sacrifice telemetry completeness for installation simplicity. Elastic observability architecture documentation describes hybrid push-pull patterns where agents push metrics and logs to centralized endpoints, accommodating firewall traversal and egress-only network policies common in enterprise settings.

Industry practices also reveal the prevalence of logging-centric observability as a pragmatic compromise when comprehensive instrumentation is infeasible. Structured logging with correlation identifiers, combined with centralized log aggregation and pattern matching, provides limited but actionable visibility without requiring application code changes or privileged host access. Jenkins and GitLab \ac{CICD} platform documentation emphasizes log-based troubleshooting workflows, reflecting the operational reality that pipeline executions are often diagnosed through console output rather than through distributed traces or granular metrics. This logging-first approach aligns with the ephemeral nature of \ac{CICD} workloads and the difficulty of instrumenting arbitrary third-party build tools. However, the literature provides minimal empirical guidance on the diagnostic effectiveness of log-only observability compared to multi-signal approaches, leaving a gap between documented best practices and evidence-based performance validation.

%----------------------------------------------------------------------------------------

\section{Identified Gaps and Implications for Constrained CI/CD Observability}
\label{sec:identified-gaps}

The synthesis of reviewed literature reveals several areas where existing research provides limited guidance for observability in privilege-constrained, heterogeneous \ac{CICD} environments. First, empirical studies evaluating non-privileged telemetry collection patterns remain scarce. While sidecar-based approaches nominally operate without host-level privileges, performance evaluations typically occur in Kubernetes clusters where orchestrator-managed injection and resource allocation mitigate operational complexity. The feasibility and overhead characteristics of sidecar telemetry in environments where container orchestration is absent or restricted—such as standalone Docker executor environments or VM-based runners—have not been systematically characterized. Similarly, the applicability of application-level instrumentation libraries to third-party packaged software, a common scenario in \ac{CICD} pipelines, receives minimal attention in the literature, which predominantly evaluates custom-built microservice applications amenable to source-level modification.

Second, the literature offers limited insight into post hoc correlation techniques suitable for ephemeral task-based workloads. Research on distributed tracing emphasizes in-band context propagation through service invocation chains, a model poorly suited to \ac{CICD} executions where observability signals must be correlated retrospectively using pipeline metadata, container identifiers, and temporal alignment of logs and metrics. Techniques for reconstructing causal relationships between resource contention events and pipeline failures, or for attributing aggregate host metrics to individual containerized build tasks, are underexplored. This gap is significant because \ac{CICD} observability fundamentally addresses a diagnostic problem—identifying which pipeline stage failed and why—rather than the performance optimization focus that dominates microservice observability research.

Third, the methodological practices documented in the literature provide insufficient guidance for evaluating observability interventions in controlled \ac{CICD} settings. Standard microservice benchmarks such as Online Boutique or DeathStarBench model synchronous request-response patterns with stateless service replicas, characteristics that diverge from \ac{CICD} workload semantics: batch job execution, stateful build caches, sequential pipeline stages with inter-stage dependencies. Metrics such as \ac{MTTD} and \ac{MTTR}, frequently cited as primary outcome measures for observability effectiveness, lack standardized operational definitions in the \ac{CICD} context. Variability in how detection events are defined (alert trigger, operator awareness, automated remediation) and how recovery is measured (pipeline restart, job completion, service restoration) complicates the interpretation and comparison of reported improvements. The absence of reproducible benchmark artifacts—representative \ac{CICD} workloads, fault injection scenarios, instrumentation templates—further hinders empirical validation of observability claims.

Finally, the literature exhibits limited engagement with operational constraints imposed by multi-tenant shared infrastructure, data residency requirements, and restricted network egress policies. Research prototypes typically assume unimpeded network connectivity to backend observability platforms, whereas enterprise \ac{CICD} environments may mandate on-premise data retention, prohibit external telemetry export, or enforce egress filtering that complicates agent-to-collector communication. The implications of such constraints for collector topology, buffering strategies, and telemetry completeness remain largely unaddressed. These gaps collectively indicate that while the observability research community has produced substantial evidence on instrumentation techniques, performance tradeoffs, and architectural patterns, the direct applicability of these findings to constrained \ac{CICD} contexts requires further investigation and adaptation.

%----------------------------------------------------------------------------------------

\section{Summary and Implications for Subsequent Research}
\label{sec:chapter2-summary}

This chapter has established the methodological foundation for systematic literature review and synthesized existing research on container observability, with particular attention to deployment assumptions, performance characteristics, and the alignment between documented practices and the operational constraints of privilege-restricted \ac{CICD} environments. The reviewed literature demonstrates that contemporary observability systems achieve rich telemetry collection through diverse instrumentation techniques—host-based kernel probes, sidecar proxies, application-level tracing libraries—each exhibiting distinct privilege requirements, performance overheads, and applicability boundaries. While kernel-level \ac{eBPF} mechanisms offer substantial efficiency advantages, their dependence on host access and recent kernel features limits transferability to heterogeneous or restricted platforms. Sidecar patterns provide deployment flexibility but introduce resource multipliers and microarchitectural contention. Application-level instrumentation enables fine-grained observability at the cost of code modification and cross-platform compatibility challenges.

The comparative analysis reveals that much of the published research presumes operational capabilities unavailable in constrained \ac{CICD} deployments: administrative host access, orchestrator-mediated resource management, and the ability to modify or instrument application code. Evaluation practices, while methodologically diverse, predominantly target long-lived microservice workloads rather than ephemeral task-based executions characteristic of \ac{CICD} pipelines. Metrics such as \ac{MTTD} and \ac{MTTR}, though frequently invoked, lack standardized definitions suitable for comparing observability interventions across heterogeneous systems. These findings underscore the necessity of adapting existing instrumentation techniques and evaluation methodologies to the specific constraints and workload characteristics of on-premise, multi-tenant \ac{CICD} infrastructure.

The identified gaps in existing research inform the design and evaluation strategy for subsequent phases of this work. The absence of empirical studies characterizing non-privileged telemetry collection in \ac{CICD} contexts motivates the development of a prototype observability system that operates within documented privilege boundaries while providing actionable diagnostic signals. The limited guidance on post hoc correlation for ephemeral workloads necessitates the investigation of metadata-based correlation techniques leveraging pipeline orchestration APIs and temporal alignment heuristics. The methodological gap in controlled \ac{CICD} evaluation frameworks justifies the establishment of reproducible benchmark scenarios with standardized fault injection and measurement procedures. Collectively, the evidence synthesized in this chapter establishes the research context and motivates the technical contributions described in subsequent chapters, bridging the gap between general-purpose observability research and the specific requirements of constrained \ac{CICD} observability.